\documentclass[11pt]{article}
\usepackage[margin=2.0cm]{geometry}
\usepackage[T1]{fontenc}
\usepackage{lmodern}
\usepackage{amsmath,amssymb,amsthm,mathtools}
\usepackage{booktabs,tabularx,longtable,array,graphicx,microtype,natbib}
\usepackage[font=small,labelfont=bf]{caption}
\usepackage[hidelinks,pdftitle={Supplementary information: Crop-ranking reversal}]{hyperref}
\newtheorem{theorem}{Theorem}
\newtheorem{proposition}{Proposition}
\newcommand{\E}{\mathbb E}
\newcommand{\R}{\mathbb R}
\newcommand{\CVaR}{\operatorname{CVaR}}
\newcommand{\ind}{\mathbf 1}
\newcommand{\DesignHash}{57e0689e388b}
\newcommand{\KansasYears}{8}
\newcommand{\ScoreCorn}{0.748}
\newcommand{\ScoreSoy}{0.756}
\newcommand{\ScoreWheat}{0.894}
\newcommand{\MarginCorn}{199.8}
\newcommand{\MarginSoy}{229.0}
\newcommand{\MarginWheat}{138.3}
\newcommand{\PrimaryCorn}{0.330}
\newcommand{\PrimarySoy}{0.422}
\newcommand{\PrimaryWheat}{0.248}
\newcommand{\PrimaryProfit}{198.5}
\newcommand{\PrimaryCVaR}{-47.0}
\newcommand{\PrimaryMinCVaR}{-59.5}
\newcommand{\PrimaryNoRiskCVaR}{-34.5}
\newcommand{\KKTPrimal}{1.42e-14}
\newcommand{\KKTStationarity}{2.84e-14}
\newcommand{\KKTComplementarity}{5.06e-14}
\newcommand{\PhaseCells}{165}
\newcommand{\PhaseReversal}{143}
\newcommand{\PhaseStrong}{96}
\newcommand{\PhaseMultiple}{2}
\newcommand{\BootstrapProbability}{0.969}
\newcommand{\BootstrapProbabilityLow}{0.892}
\newcommand{\BootstrapProbabilityHigh}{0.996}
\newcommand{\BootstrapFirstMean}{0.066}
\newcommand{\BootstrapFirstLow}{0.000}
\newcommand{\BootstrapFirstHigh}{0.598}
\newcommand{\DiversificationGaussianCVaR}{-36.5}
\newcommand{\DiversificationTailCVaR}{-45.7}
\newcommand{\DiversificationRiskCeiling}{-45.7}
\newcommand{\LowerTailCoefficient}{0.195}
\newcommand{\StateYears}{248}
\newcommand{\States}{31}
\newcommand{\RelativeYieldRate}{0.887}
\newcommand{\OperatingMarginRate}{0.625}
\newcommand{\RevenueRate}{0.613}
\newcommand{\TotalCostRate}{0.782}
\newcommand{\LaggedEffect}{-0.0042}
\newcommand{\LaggedLow}{-0.0123}
\newcommand{\LaggedHigh}{0.0056}
\newcommand{\InfoStrictCells}{25}
\newcommand{\InfoSubCells}{18}
\newcommand{\InfoZeroCells}{18}
\newcommand{\InfoBoundaryCells}{11}
\newcommand{\MaxInformationValue}{4.655}

\title{\textbf{Supplementary Information}\\
\large Crop-ranking reversal under downside risk and operational constraints}
\author{Xinyu Lu and Peimin Chen}
\date{}
\begin{document}
\maketitle
\tableofcontents

\section{Relationship to the supervisor draft}

The immutable supervisor draft establishes the research question and innovation
chain: suitability versus operational optimality; ranking reversal under
uncertain price, yield and cost; joint downside risk and diversification
failure; operational constraints; and the dependence of climate-information
value on flexibility. The reconstruction retains that chain and the section
architecture of introduction, literature, model, theory, numerical analysis,
discussion and conclusion.

Several mathematical statements in the draft could not be retained as written.
An undefined risk-adjusted margin gap did not imply an acreage ``if and only
if'' result; KKT stationarity omitted operational multipliers and
finite-scenario auxiliary variables; a scalar lower-tail coefficient did not
globally order portfolio CVaR across copula families; a universal unique
dependence threshold lacked monotonicity and crossing assumptions; and strict
information--flexibility complementarity was not unconditional. The main
paper replaces these statements with restricted exact theorems and reports
null, disconnected and multiple-optimum cases.

The previous repository manuscript centered the proposition that an ordinal
ranking does not identify a cardinal allocation. That statement is correct but
does not by itself answer the supervisor's mechanism question. In the
reconstructed manuscript it is treated as a boundary condition. The dominant
argument is again the conditional reversal frontier under the full stochastic
operational model.

\section{Complete model notation}

\begin{longtable}{p{2.2cm}p{11.8cm}}
\toprule
Symbol & Definition\\
\midrule
\endfirsthead
\toprule
Symbol & Definition\\
\midrule
\endhead
\(i=1,\ldots,n\) & crop index\\
\(\omega\) & joint price--yield--cost state\\
\(s_i\) & external pre-decision agronomic or predictive score\\
\(x_i\) & planted area or normalized land share\\
\(\widetilde P_i,\widetilde Y_i,\widetilde C_i\) & stochastic price, yield and cost\\
\(\widetilde\pi_i\) & stochastic per-acre margin
\(\widetilde P_i\widetilde Y_i-\widetilde C_i\)\\
\(\mu_i\) & expected per-acre margin\\
\(L(x)\) & portfolio loss \(-\widetilde\pi^\top x\)\\
\(A\) & available land\\
\(B,c\) & working-capital limit and crop cash coefficients\\
\(R,r\) & rotation rows and limits\\
\(K,k\) & contractual-minimum rows and requirements\\
\(G,h\) & shared-capacity rows and limits\\
\(\ell,u\) & crop-specific lower and upper area bounds\\
\(\alpha\) & CVaR confidence level\\
\(\kappa\) & loss-CVaR ceiling\\
\(\rho\) & normalized risk-tolerance path index\\
\(\mathcal X\) & operational feasible set\\
\(\mathcal S_\kappa\) & complete optimal set under ceiling \(\kappa\)\\
\(\sigma\) & deterministic selector for one reported optimum\\
\(\theta\) & parameter along a named joint-law family\\
\(\phi\) & operational-flexibility or buffering parameter\\
\bottomrule
\end{longtable}

\section{Full finite-scenario linear programme}

For scenario probabilities \(w_m>0\) summing to one, the implemented primal is
\begin{align}
\max_{x,v,q}\quad & \mu^\top x\\
\text{s.t.}\quad
&\ind^\top x\le A,\\
&c^\top x\le B,\\
&Rx\le r,\qquad -Kx\le-k,\qquad Gx\le h,\\
&\ell\le x\le u,\\
&v+\frac{1}{1-\alpha}\sum_{m=1}^M w_mq_m\le\kappa,\\
&-\pi_m^\top x-v-q_m\le0,\qquad m=1,\ldots,M,\\
&q_m\ge0,\qquad v\in\R.
\end{align}
There is no hidden penalty term and no post-solution repair of the optimized
policy. Repairs apply only to heuristic benchmark policies so that they are
compared inside the same operational set.

\subsection{KKT equations}

Let \(\eta\) be the CVaR-ceiling multiplier, \(t_m\) the multiplier for
scenario excess row \(m\), and \(z_m\) the multiplier on \(q_m\ge0\). Together
with the operational multipliers defined in the main text, stationarity is
\begin{align}
0={}&-\mu+\gamma\ind+\beta c+R^\top\rho-K^\top\chi+
G^\top\lambda-a+b-\sum_mt_m\pi_m,\\
0={}&\eta-\sum_mt_m,\\
0={}&\frac{\eta w_m}{1-\alpha}-t_m-z_m,\qquad m=1,\ldots,M.
\end{align}
All multipliers are nonnegative except the dual expression associated with the
free variable \(v\). Complementarity requires each multiplier times its
primal slack to be zero. The scenario equations allow fractional mass at the
VaR threshold and are valid for discrete scenario distributions.

\section{Proofs}

\subsection{Proof of the restricted reversal theorem}

Let \(z=x_1\) and \(x_2=A-z\). Expected profit is
\[
\mu_1z+\mu_2(A-z)=\mu_2A+(\mu_1-\mu_2)z.
\]
Because \(\mu_1>\mu_2\), this function is strictly increasing in \(z\).
Loss-CVaR is convex and continuous in the portfolio weights. Its sublevel set
intersected with the compact operational interval \(I\) is therefore a compact
interval \(\mathcal K_{\theta,\kappa}\). A strictly increasing linear
objective has the unique maximizer \(\sup\mathcal K_{\theta,\kappa}\).

The pair reverses when \(z<A-z\), equivalently \(z<A/2\). Hence reversal holds
if and only if the right endpoint of the entire risk-feasible interval lies
below \(A/2\), which is equivalent to no feasible point on
\([A/2,\overline z]\). If risk is nondecreasing on that upper half, the
minimum upper-half risk occurs at \(A/2\), so the upper half is infeasible if
and only if \(r_\theta(A/2)>\kappa\). Existence of a feasible lower-half point
ensures that this inequality indicates reversal rather than infeasibility.

If \(\theta\mapsto r_\theta(A/2)\) is continuous, strictly increasing and
crosses \(\kappa\), the intermediate-value theorem gives a unique
\(\theta^*\). If strictness is removed, an interval of thresholds can result.
If monotonicity is removed, the reversal set can be disconnected. If crossing
is removed, a threshold need not exist.

\subsection{Proof of complete optimality}

The finite-scenario problem is a bounded linear programme whenever
\(\mathcal X\) is nonempty and bounded. Linear-programming strong duality
implies that primal and dual feasibility plus equality of the objective values
are necessary and sufficient. Expanding equality of values across all
inequality rows yields stationarity and complementarity. The equations above
therefore form a complete certificate.

For the population convex formulation, write the Lagrangian using loss-CVaR
as a convex function. Under Slater's condition or another convex constraint
qualification, KKT conditions are necessary; convexity makes them sufficient.
Pairwise subtraction of acreage stationarity cancels \(\gamma\) but leaves
budget, rotation, contract, shared-resource and bound terms. No algebraic
operation on this local equality recovers an acreage ordering without the
global feasible geometry.

\subsection{Proof of named-family ordering}

For integrable random losses, CVaR is a law-invariant coherent risk measure and
is monotone under convex order. Thus
\(L_{\theta_1}(x)\preceq_{\mathrm{cx}}L_{\theta_2}(x)\) implies the stated
CVaR inequality at each \(x\). Pointwise inequality gives containment of the
two risk sublevel sets. Maximizing the same expected-profit objective over
nested sets gives the optimal-value order. Coordinate monotonicity does not
follow because the supporting face can change when the feasible set contracts.

\subsection{Proofs for information and flexibility}

Any no-information action defines a constant signal-contingent policy. If the
constant policy is feasible, optimization after observing a signal cannot
lower value. When one common action is posterior-optimal for every signal, the
constant policy also attains the informed optimum, proving zero value. If
unique posterior optima differ on events with positive probability and at
least one difference is strict, expected improvement is positive.

Increasing differences and nested action sets yield monotone posterior actions
under standard monotone comparative statics. Strictness on a positive-
probability event makes the information gain larger at the higher acreage-
flexibility level. This is the conditional complementarity result.

For buffering, at \(\phi=1\) all state margins equal \(\bar\pi\); posterior
optimization problems are identical and the common-action condition gives
zero information value. If information value is positive at \(\phi=0\) and
continuous along the contraction path, it must decline somewhere. Those
intervals are substitutive. Total optimized value may still rise with
\(\phi\).

\section{Calibration and constraint table}

\begin{table}[h]
\centering\small
\caption{Registered Kansas calibration and primary allocation.}
\begin{tabular}{lrrrr}
\toprule
Crop & score & mean margin & margin s.d. & full-model share\\
 & & \multicolumn{2}{c}{real 2024 US\$ per acre} &\\
\midrule
Corn & \ScoreCorn & \MarginCorn & 94.9 & \PrimaryCorn\\
Soybean & \ScoreSoy & \MarginSoy & 104.0 & \PrimarySoy\\
Winter wheat & \ScoreWheat & \MarginWheat & 49.6 & \PrimaryWheat\\
\bottomrule
\end{tabular}
\end{table}

\begin{table}[h]
\centering\small
\caption{Operational set. Coefficients use normalized land and capacity units.}
\begin{tabularx}{\textwidth}{lXXX}
\toprule
Object & Corn & Soybean & Winter wheat\\
\midrule
Lower bound & 0.05 & 0.10 & 0.05\\
Upper bound & 0.55 & 0.65 & 0.60\\
Rotation cap & 0.50 & --- & ---\\
Contract minimum & --- & 0.10 & ---\\
Planting-labour coefficient & 1.00 & 0.72 & 0.58\\
Harvest-equipment coefficient & 1.05 & 0.82 & 0.62\\
\bottomrule
\end{tabularx}
\end{table}

Total land is 1.0, operating budget is US\$330, planting-labour capacity is
0.82 and harvest-equipment capacity is 0.88. The full model uses
\(\alpha=0.95\), \(\rho=0.5\), 512 optimization scenarios and seed
202607270034.

\section{Policy comparison}

\begin{table}[h]
\centering\small
\caption{Registered primary policy outcomes. A larger loss-CVaR is worse.}
\begin{tabular}{lrrrrrl}
\toprule
Policy & Corn & Soy & Wheat & Profit & loss-CVaR & risk feasible\\
\midrule
Suitability proportional & 0.312 & 0.315 & 0.373 & 187.4 & -53.1 & yes\\
Winner take all, repaired & 0.300 & 0.100 & 0.600 & 166.6 & -57.3 & yes\\
Equal share & 0.333 & 0.333 & 0.333 & 190.4 & -51.4 & yes\\
Expected profit & 0.280 & 0.650 & 0.070 & 216.5 & -34.5 & no\\
Mean--variance & 0.280 & 0.650 & 0.070 & 216.5 & -34.5 & no\\
Full loss-CVaR & \PrimaryCorn & \PrimarySoy & \PrimaryWheat &
\PrimaryProfit & \PrimaryCVaR & yes\\
Minimum-CVaR endpoint & 0.220 & 0.180 & 0.600 & 169.0 & -59.5 & yes\\
\bottomrule
\end{tabular}
\end{table}

The identical primary expected-profit and mean--variance allocations are not a
coding shortcut: under the registered variance penalty and constraints, both
select the same boundary point. The diversification-failure test evaluates
that point under the Student-\(t\) tail law and compares it with the full
loss-CVaR solution.

\section{Phase-diagram summary}

\begin{longtable}{lrrrrrr}
\caption{Family-specific reversal counts on each Kendall-\(\tau\) path.}\\
\toprule
Family & \(\tau\) & reversal & strong & multiple & first reversal &
first strong\\
\midrule
\endfirsthead
\toprule
Family & \(\tau\) & reversal & strong & multiple & first reversal &
first strong\\
\midrule
\endhead
Clayton & 0.00 & 9 & 5 & 2 & 0.2 & 0.6\\
Clayton & 0.15 & 11 & 9 & 0 & 0.0 & 0.2\\
Clayton & 0.30 & 10 & 7 & 0 & 0.0 & 0.4\\
Clayton & 0.45 & 8 & 5 & 0 & 0.3 & 0.6\\
Clayton & 0.60 & 9 & 3 & 0 & 0.2 & 0.8\\
Gaussian & 0.00 & 11 & 10 & 0 & 0.0 & 0.1\\
Gaussian & 0.15 & 10 & 5 & 0 & 0.1 & 0.6\\
Gaussian & 0.30 & 10 & 8 & 0 & 0.1 & 0.3\\
Gaussian & 0.45 & 8 & 2 & 0 & 0.3 & 0.9\\
Gaussian & 0.60 & 10 & 5 & 0 & 0.1 & 0.6\\
Student-\(t\) & 0.00 & 10 & 9 & 0 & 0.1 & 0.2\\
Student-\(t\) & 0.15 & 10 & 9 & 0 & 0.1 & 0.2\\
Student-\(t\) & 0.30 & 11 & 9 & 0 & 0.0 & 0.2\\
Student-\(t\) & 0.45 & 8 & 7 & 0 & 0.3 & 0.4\\
Student-\(t\) & 0.60 & 8 & 3 & 0 & 0.3 & 0.8\\
\bottomrule
\end{longtable}

First-reversal and first-strong columns report the lowest risk-tolerance index
on the registered grid. Clayton at \(\tau=0.30\) is the only disconnected
selected-reversal path. The two multiple-optimum cells remain in all counts
and are resolved by separate possible, universal and selected classifications.

\section{Robustness and uncertainty}

\begin{table}[h]
\centering\small
\caption{One-at-a-time robustness. ``Selected'' and ``strong'' are counts.}
\begin{tabular}{lrrr}
\toprule
Dimension & cells & selected reversal & strong reversal\\
\midrule
Confidence level & 3 & 3 & 2\\
Constraint regime & 5 & 5 & 3\\
Dependence family & 3 & 3 & 2\\
Dependence parameter & 5 & 5 & 4\\
Marginal law & 3 & 3 & 3\\
Sample window & 3 & 3 & 3\\
Scenario count & 3 & 3 & 3\\
Random seed & 3 & 3 & 2\\
Solver & 3 & 3 & 3\\
\bottomrule
\end{tabular}
\end{table}

The 64 historical bootstrap replications produce mean shares 0.277 for corn,
0.527 for soybean and 0.196 for winter wheat. Their 95\% percentile intervals
are [0.125, 0.368], [0.258, 0.650] and [0.070, 0.404]. Reversal probability
is \BootstrapProbability{} with interval
[\BootstrapProbabilityLow,\BootstrapProbabilityHigh]. The first-reversal
index mean is \BootstrapFirstMean{} with interval
[\BootstrapFirstLow,\BootstrapFirstHigh]. The broad boundary interval is
consistent with an eight-year calibration and is not hidden by the large
number of simulated scenarios.

\section{Information--flexibility details}

The state split contains four low and four high relative-yield-advantage
years. Signal accuracy takes six values and flexibility takes six values on
each path, producing 72 cells. The region counts are \InfoStrictCells{}
strict-complementarity, \InfoSubCells{} substitution, \InfoZeroCells{} zero
and \InfoBoundaryCells{} weak or boundary. Every signal-contingent problem is
feasible and the minimum computed value of information differs from zero only
at floating-point scale.

At perfect signal accuracy, post-signal acreage reallocation raises
information value monotonically from zero at no adjustable acreage to
\(\$\,\MaxInformationValue\) at full adjustment. Under shock buffering, the
corresponding information value declines toward zero as the state margins are
contracted to a common vector. This comparison is about interaction; buffering
can still increase the no-information optimized value.

\section{External panel and identification boundary}

\begin{table}[h]
\centering\small
\caption{Top-rank disagreement in \StateYears{} state-years.}
\begin{tabular}{lrrr}
\toprule
Ranking definition & estimate & cluster 95\% low & cluster 95\% high\\
\midrule
Relative yield & \RelativeYieldRate & 0.802 & 0.954\\
Standardized revenue & \RevenueRate & 0.447 & 0.767\\
Operating margin & \OperatingMarginRate & 0.552 & 0.693\\
Total-cost margin & \TotalCostRate & 0.701 & 0.857\\
\bottomrule
\end{tabular}
\end{table}

The leakage-free relative-yield specification uses 651 crop-transition rows.
Its prior-leader effect is \LaggedEffect{} with interval
[\LaggedLow,\LaggedHigh]. The estimand is descriptive because the panel does
not observe farm objectives, local prices or costs, credit, rotation history,
contracts, risk preferences or received forecasts. State acreage is never
substituted for an optimizer output.

\section{Data lineage and reproducibility}

The reconstruction begins with registered raw USDA NASS, USDA ERS and BLS
files. The canonical state-crop panel is joined to national prices, costs and
CPI-U through explicit crop and year keys. The Issue \#34 script filters the
frozen Kansas calibration years, constructs the genuine score and real
margins, solves all policies, phase cells, robustness and bootstrap
replications, and writes a machine-readable manifest.

The required commands are:
\begin{verbatim}
uv sync
uv run python scripts/run_issue34_reconstruction.py
(cd reconstruction/issue34/outputs && \
  shasum -a 256 -c SHA256SUMS.txt)
uv run python scripts/render_issue34_manuscript_numbers.py
uv run python scripts/make_issue34_figures.py
uv run pytest -q
\end{verbatim}

The plotting script reads only registered output tables. Each figure is
exported as editable SVG and PDF, 300-dpi PNG and 600-dpi TIFF. The manuscript
and Supplementary Information compile from the repository root. The release
archive contains the source registry, data dictionary, analysis manifest,
checksums, tests, source data and rendered documents.

\section{Limitations}

\begin{enumerate}
\item The Kansas calibration has eight effective annual observations. It
cannot estimate extreme-tail dependence reliably.
\item State yield combined with national price and cost is not a farm-level
profit series.
\item Copula parameters and shared-capacity limits are registered stress
values, not causal estimates.
\item The model is single-period and linear. It omits fixed charges,
endogenous prices, spatial spillovers and explicit multi-year rotation state.
\item Historical bootstrap replications quantify sensitivity to the short
calibration record but do not create independent empirical observations.
\item The information state is a transparent annual proxy, not a deployed
forecast product.
\item The external panel is aggregate and descriptive; it cannot identify the
full model's active mechanism or welfare.
\item Author order, affiliations and contribution wording require final
confirmation before journal submission.
\end{enumerate}

\end{document}
